% settings.tex — packages, theme, colors, macros for the SEAI Minesweeper RL slides.
% Included from slides.tex via % settings.tex — packages, theme, colors, macros for the SEAI Minesweeper RL slides.
% Included from slides.tex via % settings.tex — packages, theme, colors, macros for the SEAI Minesweeper RL slides.
% Included from slides.tex via % settings.tex — packages, theme, colors, macros for the SEAI Minesweeper RL slides.
% Included from slides.tex via \input{settings} AFTER \documentclass.

% ---------------------------------------------------------------------------
% Theme
% ---------------------------------------------------------------------------
\usetheme{metropolis}

% ---------------------------------------------------------------------------
% Encoding / fonts / language
% ---------------------------------------------------------------------------
\usepackage[utf8]{inputenc}
\usepackage[T1]{fontenc}
\usepackage[english]{babel}
\usepackage{csquotes} % biblatex recommends this alongside babel
\usepackage{microtype}
\usepackage{lmodern}

% ---------------------------------------------------------------------------
% Math
% ---------------------------------------------------------------------------
\usepackage{amsmath}
\usepackage{amssymb}
\usepackage{mathtools}
\usepackage{bm}

% ---------------------------------------------------------------------------
% Graphics / figures / plots / diagrams
% ---------------------------------------------------------------------------
\usepackage{graphicx}
\usepackage{adjustbox} % valign=c: vertically center images against text on the same line
\usepackage{caption}
\usepackage{subcaption}
\usepackage{tikz}
\usetikzlibrary{arrows.meta, positioning, calc, shapes.geometric, fit, backgrounds, decorations.pathreplacing}
\usepackage{pgfplots}
\pgfplotsset{compat=1.18}

% ---------------------------------------------------------------------------
% Tables
% ---------------------------------------------------------------------------
\usepackage{booktabs}
\usepackage{tabularx}
\usepackage{siunitx}

% ---------------------------------------------------------------------------
% Algorithms / code listings
% ---------------------------------------------------------------------------
\usepackage[ruled, vlined, linesnumbered]{algorithm2e}
\usepackage{listings}

% ---------------------------------------------------------------------------
% Callout boxes
% ---------------------------------------------------------------------------
\usepackage[most]{tcolorbox}
\tcbuselibrary{skins}

% ---------------------------------------------------------------------------
% Icons, hyperlinks
% ---------------------------------------------------------------------------
\usepackage{fontawesome5}
\usepackage{hyperref}

% ---------------------------------------------------------------------------
% Bibliography (biblatex + biber)
% ---------------------------------------------------------------------------
\usepackage[
  backend=biber,
  style=authoryear,
  sorting=nyt,
  maxbibnames=3,
  giveninits=true,
]{biblatex}
\addbibresource{references.bib}

% ---------------------------------------------------------------------------
% Colors — accent taken from the UniPi seal (sampled: #00518b)
% ---------------------------------------------------------------------------
\definecolor{unipiblue}{HTML}{00518B}
\definecolor{unipiblueDark}{HTML}{00365C}
\definecolor{unipiblueLight}{HTML}{4D89B3}
\definecolor{accentgray}{HTML}{3B3B3B}
\definecolor{lightgray}{HTML}{F2F2F2}
\definecolor{goodgreen}{HTML}{2E7D32}
\definecolor{badred}{HTML}{C62828}

\setbeamercolor{palette primary}{bg=unipiblue, fg=white}
\setbeamercolor{frametitle}{bg=unipiblue, fg=white}
\setbeamercolor{progress bar}{fg=unipiblue, bg=lightgray}
\setbeamercolor{alerted text}{fg=badred}
\setbeamercolor{example text}{fg=goodgreen}
\setbeamercolor{normal text}{fg=accentgray}

% ---------------------------------------------------------------------------
% Metropolis / general Beamer tweaks
% ---------------------------------------------------------------------------
\metroset{
  progressbar=frametitle,
  numbering=fraction,
  block=fill,
  sectionpage=progressbar,
}
\setbeamertemplate{navigation symbols}{}
\setbeamertemplate{itemize items}[circle]
\setbeamertemplate{caption}[numbered]
\captionsetup{font=scriptsize, labelfont={bf,color=unipiblue}}

% ---------------------------------------------------------------------------
% Custom macros — math / RL notation
% ---------------------------------------------------------------------------
\DeclareMathOperator*{\argmax}{arg\,max}
\DeclareMathOperator*{\argmin}{arg\,min}
\newcommand{\E}[1]{\mathbb{E}\left[#1\right]}
\newcommand{\Prob}[1]{\mathbb{P}\left(#1\right)}
\newcommand{\states}{\mathcal{S}}
\newcommand{\actions}{\mathcal{A}}
\newcommand{\reward}{\mathcal{R}}
\newcommand{\policy}{\pi}
\newcommand{\qfunc}[2]{Q\left(#1, #2\right)}
\newcommand{\qstar}[2]{Q^{*}\left(#1, #2\right)}
\newcommand{\vfunc}[1]{V\left(#1\right)}
\newcommand{\replaybuffer}{\mathcal{D}}
\newcommand{\dqn}{\textsc{dqn}}
\newcommand{\per}{\textsc{per}}
\newcommand{\ddqn}{\textsc{c51}}

% ---------------------------------------------------------------------------
% Custom environments / commands — slide-building helpers
% ---------------------------------------------------------------------------

% Colored callout box: \highlightbox{title}{body}
\newtcolorbox{highlightbox}[2][]{
  colback=lightgray, colframe=unipiblue, coltitle=white,
  fonttitle=\bfseries, title=#2, boxrule=0.5pt, arc=2pt, #1
}

% Full-bleed figure with caption, sized to a fraction of textwidth:
% \fullfigure{width-fraction}{path}{caption}{label}
\newcommand{\fullfigure}[4]{%
  \begin{figure}
    \centering
    \includegraphics[width=#1\textwidth]{#2}
    \caption{#3}
    \label{#4}
  \end{figure}
}

% Two-column slide split: \twocol{left content}{right content}
\newcommand{\twocol}[2]{%
  \begin{columns}[T, totalwidth=\textwidth]
    \begin{column}{0.48\textwidth}
      #1
    \end{column}
    \begin{column}{0.48\textwidth}
      #2
    \end{column}
  \end{columns}
}

% TikZ style shorthands for MDP / architecture diagrams
\tikzset{
state/.style={circle, draw=unipiblue, thick, minimum size=1cm, fill=lightgray},
action/.style={rectangle, draw=accentgray, thick, minimum size=0.8cm, rounded corners},
flowarrow/.style={-{Latex[length=2mm]}, thick, unipiblue},
}
 AFTER \documentclass.

% ---------------------------------------------------------------------------
% Theme
% ---------------------------------------------------------------------------
\usetheme{metropolis}

% ---------------------------------------------------------------------------
% Encoding / fonts / language
% ---------------------------------------------------------------------------
\usepackage[utf8]{inputenc}
\usepackage[T1]{fontenc}
\usepackage[english]{babel}
\usepackage{csquotes} % biblatex recommends this alongside babel
\usepackage{microtype}
\usepackage{lmodern}

% ---------------------------------------------------------------------------
% Math
% ---------------------------------------------------------------------------
\usepackage{amsmath}
\usepackage{amssymb}
\usepackage{mathtools}
\usepackage{bm}

% ---------------------------------------------------------------------------
% Graphics / figures / plots / diagrams
% ---------------------------------------------------------------------------
\usepackage{graphicx}
\usepackage{adjustbox} % valign=c: vertically center images against text on the same line
\usepackage{caption}
\usepackage{subcaption}
\usepackage{tikz}
\usetikzlibrary{arrows.meta, positioning, calc, shapes.geometric, fit, backgrounds, decorations.pathreplacing}
\usepackage{pgfplots}
\pgfplotsset{compat=1.18}

% ---------------------------------------------------------------------------
% Tables
% ---------------------------------------------------------------------------
\usepackage{booktabs}
\usepackage{tabularx}
\usepackage{siunitx}

% ---------------------------------------------------------------------------
% Algorithms / code listings
% ---------------------------------------------------------------------------
\usepackage[ruled, vlined, linesnumbered]{algorithm2e}
\usepackage{listings}

% ---------------------------------------------------------------------------
% Callout boxes
% ---------------------------------------------------------------------------
\usepackage[most]{tcolorbox}
\tcbuselibrary{skins}

% ---------------------------------------------------------------------------
% Icons, hyperlinks
% ---------------------------------------------------------------------------
\usepackage{fontawesome5}
\usepackage{hyperref}

% ---------------------------------------------------------------------------
% Bibliography (biblatex + biber)
% ---------------------------------------------------------------------------
\usepackage[
  backend=biber,
  style=authoryear,
  sorting=nyt,
  maxbibnames=3,
  giveninits=true,
]{biblatex}
\addbibresource{references.bib}

% ---------------------------------------------------------------------------
% Colors — accent taken from the UniPi seal (sampled: #00518b)
% ---------------------------------------------------------------------------
\definecolor{unipiblue}{HTML}{00518B}
\definecolor{unipiblueDark}{HTML}{00365C}
\definecolor{unipiblueLight}{HTML}{4D89B3}
\definecolor{accentgray}{HTML}{3B3B3B}
\definecolor{lightgray}{HTML}{F2F2F2}
\definecolor{goodgreen}{HTML}{2E7D32}
\definecolor{badred}{HTML}{C62828}

\setbeamercolor{palette primary}{bg=unipiblue, fg=white}
\setbeamercolor{frametitle}{bg=unipiblue, fg=white}
\setbeamercolor{progress bar}{fg=unipiblue, bg=lightgray}
\setbeamercolor{alerted text}{fg=badred}
\setbeamercolor{example text}{fg=goodgreen}
\setbeamercolor{normal text}{fg=accentgray}

% ---------------------------------------------------------------------------
% Metropolis / general Beamer tweaks
% ---------------------------------------------------------------------------
\metroset{
  progressbar=frametitle,
  numbering=fraction,
  block=fill,
  sectionpage=progressbar,
}
\setbeamertemplate{navigation symbols}{}
\setbeamertemplate{itemize items}[circle]
\setbeamertemplate{caption}[numbered]
\captionsetup{font=scriptsize, labelfont={bf,color=unipiblue}}

% ---------------------------------------------------------------------------
% Custom macros — math / RL notation
% ---------------------------------------------------------------------------
\DeclareMathOperator*{\argmax}{arg\,max}
\DeclareMathOperator*{\argmin}{arg\,min}
\newcommand{\E}[1]{\mathbb{E}\left[#1\right]}
\newcommand{\Prob}[1]{\mathbb{P}\left(#1\right)}
\newcommand{\states}{\mathcal{S}}
\newcommand{\actions}{\mathcal{A}}
\newcommand{\reward}{\mathcal{R}}
\newcommand{\policy}{\pi}
\newcommand{\qfunc}[2]{Q\left(#1, #2\right)}
\newcommand{\qstar}[2]{Q^{*}\left(#1, #2\right)}
\newcommand{\vfunc}[1]{V\left(#1\right)}
\newcommand{\replaybuffer}{\mathcal{D}}
\newcommand{\dqn}{\textsc{dqn}}
\newcommand{\per}{\textsc{per}}
\newcommand{\ddqn}{\textsc{c51}}

% ---------------------------------------------------------------------------
% Custom environments / commands — slide-building helpers
% ---------------------------------------------------------------------------

% Colored callout box: \highlightbox{title}{body}
\newtcolorbox{highlightbox}[2][]{
  colback=lightgray, colframe=unipiblue, coltitle=white,
  fonttitle=\bfseries, title=#2, boxrule=0.5pt, arc=2pt, #1
}

% Full-bleed figure with caption, sized to a fraction of textwidth:
% \fullfigure{width-fraction}{path}{caption}{label}
\newcommand{\fullfigure}[4]{%
  \begin{figure}
    \centering
    \includegraphics[width=#1\textwidth]{#2}
    \caption{#3}
    \label{#4}
  \end{figure}
}

% Two-column slide split: \twocol{left content}{right content}
\newcommand{\twocol}[2]{%
  \begin{columns}[T, totalwidth=\textwidth]
    \begin{column}{0.48\textwidth}
      #1
    \end{column}
    \begin{column}{0.48\textwidth}
      #2
    \end{column}
  \end{columns}
}

% TikZ style shorthands for MDP / architecture diagrams
\tikzset{
state/.style={circle, draw=unipiblue, thick, minimum size=1cm, fill=lightgray},
action/.style={rectangle, draw=accentgray, thick, minimum size=0.8cm, rounded corners},
flowarrow/.style={-{Latex[length=2mm]}, thick, unipiblue},
}
 AFTER \documentclass.

% ---------------------------------------------------------------------------
% Theme
% ---------------------------------------------------------------------------
\usetheme{metropolis}

% ---------------------------------------------------------------------------
% Encoding / fonts / language
% ---------------------------------------------------------------------------
\usepackage[utf8]{inputenc}
\usepackage[T1]{fontenc}
\usepackage[english]{babel}
\usepackage{csquotes} % biblatex recommends this alongside babel
\usepackage{microtype}
\usepackage{lmodern}

% ---------------------------------------------------------------------------
% Math
% ---------------------------------------------------------------------------
\usepackage{amsmath}
\usepackage{amssymb}
\usepackage{mathtools}
\usepackage{bm}

% ---------------------------------------------------------------------------
% Graphics / figures / plots / diagrams
% ---------------------------------------------------------------------------
\usepackage{graphicx}
\usepackage{adjustbox} % valign=c: vertically center images against text on the same line
\usepackage{caption}
\usepackage{subcaption}
\usepackage{tikz}
\usetikzlibrary{arrows.meta, positioning, calc, shapes.geometric, fit, backgrounds, decorations.pathreplacing}
\usepackage{pgfplots}
\pgfplotsset{compat=1.18}

% ---------------------------------------------------------------------------
% Tables
% ---------------------------------------------------------------------------
\usepackage{booktabs}
\usepackage{tabularx}
\usepackage{siunitx}

% ---------------------------------------------------------------------------
% Algorithms / code listings
% ---------------------------------------------------------------------------
\usepackage[ruled, vlined, linesnumbered]{algorithm2e}
\usepackage{listings}

% ---------------------------------------------------------------------------
% Callout boxes
% ---------------------------------------------------------------------------
\usepackage[most]{tcolorbox}
\tcbuselibrary{skins}

% ---------------------------------------------------------------------------
% Icons, hyperlinks
% ---------------------------------------------------------------------------
\usepackage{fontawesome5}
\usepackage{hyperref}

% ---------------------------------------------------------------------------
% Bibliography (biblatex + biber)
% ---------------------------------------------------------------------------
\usepackage[
  backend=biber,
  style=authoryear,
  sorting=nyt,
  maxbibnames=3,
  giveninits=true,
]{biblatex}
\addbibresource{references.bib}

% ---------------------------------------------------------------------------
% Colors — accent taken from the UniPi seal (sampled: #00518b)
% ---------------------------------------------------------------------------
\definecolor{unipiblue}{HTML}{00518B}
\definecolor{unipiblueDark}{HTML}{00365C}
\definecolor{unipiblueLight}{HTML}{4D89B3}
\definecolor{accentgray}{HTML}{3B3B3B}
\definecolor{lightgray}{HTML}{F2F2F2}
\definecolor{goodgreen}{HTML}{2E7D32}
\definecolor{badred}{HTML}{C62828}

\setbeamercolor{palette primary}{bg=unipiblue, fg=white}
\setbeamercolor{frametitle}{bg=unipiblue, fg=white}
\setbeamercolor{progress bar}{fg=unipiblue, bg=lightgray}
\setbeamercolor{alerted text}{fg=badred}
\setbeamercolor{example text}{fg=goodgreen}
\setbeamercolor{normal text}{fg=accentgray}

% ---------------------------------------------------------------------------
% Metropolis / general Beamer tweaks
% ---------------------------------------------------------------------------
\metroset{
  progressbar=frametitle,
  numbering=fraction,
  block=fill,
  sectionpage=progressbar,
}
\setbeamertemplate{navigation symbols}{}
\setbeamertemplate{itemize items}[circle]
\setbeamertemplate{caption}[numbered]
\captionsetup{font=scriptsize, labelfont={bf,color=unipiblue}}

% ---------------------------------------------------------------------------
% Custom macros — math / RL notation
% ---------------------------------------------------------------------------
\DeclareMathOperator*{\argmax}{arg\,max}
\DeclareMathOperator*{\argmin}{arg\,min}
\newcommand{\E}[1]{\mathbb{E}\left[#1\right]}
\newcommand{\Prob}[1]{\mathbb{P}\left(#1\right)}
\newcommand{\states}{\mathcal{S}}
\newcommand{\actions}{\mathcal{A}}
\newcommand{\reward}{\mathcal{R}}
\newcommand{\policy}{\pi}
\newcommand{\qfunc}[2]{Q\left(#1, #2\right)}
\newcommand{\qstar}[2]{Q^{*}\left(#1, #2\right)}
\newcommand{\vfunc}[1]{V\left(#1\right)}
\newcommand{\replaybuffer}{\mathcal{D}}
\newcommand{\dqn}{\textsc{dqn}}
\newcommand{\per}{\textsc{per}}
\newcommand{\ddqn}{\textsc{c51}}

% ---------------------------------------------------------------------------
% Custom environments / commands — slide-building helpers
% ---------------------------------------------------------------------------

% Colored callout box: \highlightbox{title}{body}
\newtcolorbox{highlightbox}[2][]{
  colback=lightgray, colframe=unipiblue, coltitle=white,
  fonttitle=\bfseries, title=#2, boxrule=0.5pt, arc=2pt, #1
}

% Full-bleed figure with caption, sized to a fraction of textwidth:
% \fullfigure{width-fraction}{path}{caption}{label}
\newcommand{\fullfigure}[4]{%
  \begin{figure}
    \centering
    \includegraphics[width=#1\textwidth]{#2}
    \caption{#3}
    \label{#4}
  \end{figure}
}

% Two-column slide split: \twocol{left content}{right content}
\newcommand{\twocol}[2]{%
  \begin{columns}[T, totalwidth=\textwidth]
    \begin{column}{0.48\textwidth}
      #1
    \end{column}
    \begin{column}{0.48\textwidth}
      #2
    \end{column}
  \end{columns}
}

% TikZ style shorthands for MDP / architecture diagrams
\tikzset{
state/.style={circle, draw=unipiblue, thick, minimum size=1cm, fill=lightgray},
action/.style={rectangle, draw=accentgray, thick, minimum size=0.8cm, rounded corners},
flowarrow/.style={-{Latex[length=2mm]}, thick, unipiblue},
}
 AFTER \documentclass.

% ---------------------------------------------------------------------------
% Theme
% ---------------------------------------------------------------------------
\usetheme{metropolis}

% ---------------------------------------------------------------------------
% Encoding / fonts / language
% ---------------------------------------------------------------------------
\usepackage[utf8]{inputenc}
\usepackage[T1]{fontenc}
\usepackage[english]{babel}
\usepackage{csquotes} % biblatex recommends this alongside babel
\usepackage{microtype}
\usepackage{lmodern}

% ---------------------------------------------------------------------------
% Math
% ---------------------------------------------------------------------------
\usepackage{amsmath}
\usepackage{amssymb}
\usepackage{mathtools}
\usepackage{bm}

% ---------------------------------------------------------------------------
% Graphics / figures / plots / diagrams
% ---------------------------------------------------------------------------
\usepackage{graphicx}
\usepackage{adjustbox} % valign=c: vertically center images against text on the same line
\usepackage{caption}
\usepackage{subcaption}
\usepackage{tikz}
\usetikzlibrary{arrows.meta, positioning, calc, shapes.geometric, fit, backgrounds, decorations.pathreplacing}
\usepackage{pgfplots}
\pgfplotsset{compat=1.18}

% ---------------------------------------------------------------------------
% Tables
% ---------------------------------------------------------------------------
\usepackage{booktabs}
\usepackage{tabularx}
\usepackage{siunitx}

% ---------------------------------------------------------------------------
% Algorithms / code listings
% ---------------------------------------------------------------------------
\usepackage[ruled, vlined, linesnumbered]{algorithm2e}
\usepackage{listings}

% ---------------------------------------------------------------------------
% Callout boxes
% ---------------------------------------------------------------------------
\usepackage[most]{tcolorbox}
\tcbuselibrary{skins}

% ---------------------------------------------------------------------------
% Icons, hyperlinks
% ---------------------------------------------------------------------------
\usepackage{fontawesome5}
\usepackage{hyperref}

% ---------------------------------------------------------------------------
% Bibliography (biblatex + biber)
% ---------------------------------------------------------------------------
\usepackage[
  backend=biber,
  style=authoryear,
  sorting=nyt,
  maxbibnames=3,
  giveninits=true,
]{biblatex}
\addbibresource{references.bib}

% ---------------------------------------------------------------------------
% Colors — accent taken from the UniPi seal (sampled: #00518b)
% ---------------------------------------------------------------------------
\definecolor{unipiblue}{HTML}{00518B}
\definecolor{unipiblueDark}{HTML}{00365C}
\definecolor{unipiblueLight}{HTML}{4D89B3}
\definecolor{accentgray}{HTML}{3B3B3B}
\definecolor{lightgray}{HTML}{F2F2F2}
\definecolor{goodgreen}{HTML}{2E7D32}
\definecolor{badred}{HTML}{C62828}

\setbeamercolor{palette primary}{bg=unipiblue, fg=white}
\setbeamercolor{frametitle}{bg=unipiblue, fg=white}
\setbeamercolor{progress bar}{fg=unipiblue, bg=lightgray}
\setbeamercolor{alerted text}{fg=badred}
\setbeamercolor{example text}{fg=goodgreen}
\setbeamercolor{normal text}{fg=accentgray}

% ---------------------------------------------------------------------------
% Metropolis / general Beamer tweaks
% ---------------------------------------------------------------------------
\metroset{
  progressbar=frametitle,
  numbering=fraction,
  block=fill,
  sectionpage=progressbar,
}
\setbeamertemplate{navigation symbols}{}
\setbeamertemplate{itemize items}[circle]
\setbeamertemplate{caption}[numbered]
\captionsetup{font=scriptsize, labelfont={bf,color=unipiblue}}

% ---------------------------------------------------------------------------
% Custom macros — math / RL notation
% ---------------------------------------------------------------------------
\DeclareMathOperator*{\argmax}{arg\,max}
\DeclareMathOperator*{\argmin}{arg\,min}
\newcommand{\E}[1]{\mathbb{E}\left[#1\right]}
\newcommand{\Prob}[1]{\mathbb{P}\left(#1\right)}
\newcommand{\states}{\mathcal{S}}
\newcommand{\actions}{\mathcal{A}}
\newcommand{\reward}{\mathcal{R}}
\newcommand{\policy}{\pi}
\newcommand{\qfunc}[2]{Q\left(#1, #2\right)}
\newcommand{\qstar}[2]{Q^{*}\left(#1, #2\right)}
\newcommand{\vfunc}[1]{V\left(#1\right)}
\newcommand{\replaybuffer}{\mathcal{D}}
\newcommand{\dqn}{\textsc{dqn}}
\newcommand{\per}{\textsc{per}}
\newcommand{\ddqn}{\textsc{c51}}

% ---------------------------------------------------------------------------
% Custom environments / commands — slide-building helpers
% ---------------------------------------------------------------------------

% Colored callout box: \highlightbox{title}{body}
\newtcolorbox{highlightbox}[2][]{
  colback=lightgray, colframe=unipiblue, coltitle=white,
  fonttitle=\bfseries, title=#2, boxrule=0.5pt, arc=2pt, #1
}

% Full-bleed figure with caption, sized to a fraction of textwidth:
% \fullfigure{width-fraction}{path}{caption}{label}
\newcommand{\fullfigure}[4]{%
  \begin{figure}
    \centering
    \includegraphics[width=#1\textwidth]{#2}
    \caption{#3}
    \label{#4}
  \end{figure}
}

% Two-column slide split: \twocol{left content}{right content}
\newcommand{\twocol}[2]{%
  \begin{columns}[T, totalwidth=\textwidth]
    \begin{column}{0.48\textwidth}
      #1
    \end{column}
    \begin{column}{0.48\textwidth}
      #2
    \end{column}
  \end{columns}
}

% TikZ style shorthands for MDP / architecture diagrams
\tikzset{
state/.style={circle, draw=unipiblue, thick, minimum size=1cm, fill=lightgray},
action/.style={rectangle, draw=accentgray, thick, minimum size=0.8cm, rounded corners},
flowarrow/.style={-{Latex[length=2mm]}, thick, unipiblue},
}
